\chapter{Considerações finais}
\label{sec:conclusao}

A proibição da pulverização aérea no Ceará é uma política de saúde com quase
seis anos de vigência plena, um vizinho que não a adotou e nenhuma avaliação.
Este projeto propõe avaliá-la com um desenho de fronteira na Chapada do Apodi,
organizado pela teoria da mudança da lei: os elos observáveis da cadeia têm cada
um a sua medida, e o efeito estimado só é atribuído à hipótese da política
depois de examinados os caminhos concorrentes. Da avaliação do desenho sai uma
recomendação que não depende dos resultados: a política carece de
monitoramento, e os registros de aplicação e de receituário permitiriam
instituí-lo.

O projeto tem limitações que convém declarar. A exposição da gestante não é
observada, e todos os efeitos são de intenção de tratar no nível do município.
Há dois ou três municípios tratados no núcleo, e os métodos de inferência para
esse caso têm baixo poder. O primeiro estágio depende de pedidos de acesso à
informação, que podem ser negados ou devolver registros incompletos. A linha de
base do método de aplicação é de 2006. Um choque estadual simultâneo à lei só é
tratado pelas especificações secundárias. E a análise não é inteiramente cega: a
série estadual de produção e os números de dois municípios tratados são
públicos e foram vistos, e uma versão anterior do projeto, com outro desenho e
só com municípios cearenses, estimou efeitos sobre o peso ao nascer depois de
2019.
